\section{A marked family in a fixed progression}
\label{sec:marked-family}

We now prepare the weighted family used to prove
Proposition~\ref{prop:reservoir}.  Fix an odd squarefree positive integer
\(M\) and a residue \(a\pmod M\) satisfying
\begin{equation}\label{eq:admissible-progression}
 \gcd(a(a-1),M)=1.
\end{equation}
The case \(M=1\) means the trivial progression.  Throughout this section,
\(x\) is real and tends to infinity, and
\begin{equation}\label{eq:marked-parameters}
 L=\log x,\qquad W=\exp(L^{0.24}),\qquad
 V(y)=\prod_{p\le y}\left(1-\frac1p\right)\quad(y\ge1).
\end{equation}
Products indexed by \(p\) are over primes.  Write \(P^-(h)\) for the least
prime factor of \(h>1\), and put \(P^-(1)=\infty\).  An integer with
\(P^-(h)>W\) is called \(W\)-rough.  Every parameter below, apart from an
explicitly displayed dependence on \(x\), is fixed before \(x\) grows.

For an integer \(K\ge1\) and fixed exponents
\(0.1<a_1<\cdots<a_K<0.2\), define the prime groups
\begin{equation}\label{eq:marked-groups}
 \mathcal P_i=\{p\text{ prime}:\exp(L^{a_i})\le p\le\exp(2L^{a_i})\},
 \qquad
 V_i=\sum_{p\in\mathcal P_i}\frac1p.
\end{equation}
These groups are disjoint for sufficiently large \(x\), and Mertens'
theorem gives \(V_i=\log2+o(1)\).  We shall choose the sieve parameters
first, then a sufficiently large \(K\), and then any admissible fixed
exponents \(a_i\).  Thresholds for \(x\) may depend on all these choices.

For a positive integer \(h\), let
\begin{equation}\label{eq:complete-group-part}
 h_{\mathcal P}=\prod_{p\in\bigcup_i\mathcal P_i}p^{v_p(h)},
 \qquad
 \omega_i(h)=\sum_{p\in\mathcal P_i}\one_{p\mid h}.
\end{equation}
Thus \(h_{\mathcal P}\) includes every power of every group prime dividing
\(h\).  Let \(\mathcal G\) be the set of positive integers supported on
the group primes, including \(1\); we call these the group integers.
The marks of \cite[Section~10]{PrimitiveRoot2}, with mark parameter
\(1/2\), are
\begin{equation}\label{eq:marked-weight}
 \mathcal W(h)=2^{K-\sum_i\omega_i(h)}
                    \prod_{i=1}^K\frac{\omega_i(h)}{V_i}.
\end{equation}
They vanish unless every group supplies a prime divisor.  Since
\(j2^{1-j}\le1\) for every integer \(j\ge1\), for sufficiently large
\(x\) we have
\begin{equation}\label{eq:mark-bound}
 0\le\mathcal W(h)\le B_K,\qquad B_K=\left(\frac2{\log2}\right)^K,
 \quad h\ge1.
\end{equation}
The constant \(B_K\) does not depend on the exponents; the threshold may.
A function \(F\) on the positive integers satisfies \(F(ph)=F(h)\) for
every group prime \(p\) and every \(h\ge1\) if and only if it depends only
on \(h/h_{\mathcal P}\).  This equivalence includes the case \(p\mid h\),
so the invariance removes all powers of a group prime.
In the application, \(F\) will select \(h/h_{\mathcal P}=4Q\) with
\(Q>x^{0.9}\) prime, so the complete group part supplies the factor \(r\)
in \(d-1=4rQ\).

The rough-factor density used below is, for \(w>\gamma>0\),
\begin{equation}\label{eq:rough-density}
 D_\gamma(w)=\frac1w+
 \sum_{j\ge2}\frac1{j!}
 \int_{\substack{t_i\ge\gamma\ (1\le i<j)\\
                  \sum_{i<j}t_i\le w-\gamma}}
 \frac1{w-\sum_{i<j}t_i}\prod_{i<j}\frac{\dd t_i}{t_i}.
\end{equation}
This is the normalization in \cite[Section~10]{PrimitiveRoot2}.  Its
\(j\)-th term is zero unless \(w\ge j\gamma\), so the sum is locally
finite.  It is nonnegative and jointly continuous on \(w>\gamma>0\).
For completeness, in a term with \(j\ge2\) put \(s=w-j\gamma\).  When
\(s\ge0\), the substitution \(t_i=\gamma+s y_i\) writes that term as
\[
 \frac{s^{j-1}}{j!}
 \int_{\substack{y_i\ge0\\\sum_{i<j}y_i\le1}}
 \frac{\prod_{i<j}\dd y_i}
 {\prod_{i<j}(\gamma+s y_i)
       \bigl(\gamma+s(1-\sum_{i<j}y_i)\bigr)}.
\]
On any compact subset of \(w>\gamma>0\) the denominators stay bounded
away from zero.  Extending this expression by zero for \(s<0\) proves
continuity also at \(w=j\gamma\).  Local finiteness completes the claim.

All Dirichlet characters in the following input, including principal
characters, are extended by zero on nonunits.  All frequencies are real.
Part (i) of the next input transfers scalar character cancellation to the
marked bilinear sum.  Part (ii) supplies that cancellation for the
difference of the two roughness tests, with
\(D_\gamma(\log m/L)/(LV(W))\) as the comparison factor.
These are the coefficient criterion of
\cite[Lemma~10.2]{PrimitiveRoot2}, specialized to the marks above, and the
cancellation verified in the proof of
\cite[Proposition~10.3]{PrimitiveRoot2}.

\begin{hypothesis}[Marked Type II estimate]\label{input:marked}
\emph{(i) Coefficient criterion.}
Fix \(\delta,C,D_*>0\).  Suppose \(H_m,H_n\ge x^\delta\) and
\(X=H_mH_n\asymp x\), with fixed comparison constants.  Let
\(F:\mathbb Z_{>0}\to\mathbb C\) satisfy \(|F(h)|\le1\) and
\(F(ph)=F(h)\) for every group prime \(p\) and every \(h\ge1\).
Let \((\alpha_m)\) and \((\beta_n)\) be scalar sequences supported on
\(W\)-rough integers in \([H_m,2H_m]\) and \([H_n,2H_n]\), respectively,
with \(\alpha\) supported on an arbitrary interval
\(I\subseteq[H_m,2H_m]\), and with
\(|\alpha_m|,|\beta_n|\le L^C\).  Suppose that for every fixed
\(A_0,B,A>0\), every Dirichlet character \(\chi\) modulo
\(k\le L^{A_0}\), and every \(|t|\le2XL^B\),
\begin{equation}\label{eq:marked-cancellation}
 \left|\frac1{H_m}\sum_m\alpha_m\chi(m)m^{it}\right|\ll_A L^{-A},
\end{equation}
uniformly in the permitted interval, character, and frequency.  The
constant and threshold in this hypothesis may depend on \(A_0,B,A\),
the fixed setup, and all further fixed parameters defining the
coefficients.  Then
\begin{equation}\label{eq:marked-input-bound}
 \left|\sum_{m,n}\alpha_m\beta_n F(mn-1)\mathcal W(mn-1)\right|
       \ll XL^{-D_*}
\end{equation}
when \(K\) is sufficiently large in terms of the fixed data, independently
of the particular \(a_i\).  Its implicit constant and threshold may depend
on all fixed data, including the common constants and thresholds in
\eqref{eq:marked-cancellation} and the \(a_i\).  With those data fixed, the
bound is uniform over the permitted \(F\), dyads, intervals, and
coefficients.

\emph{(ii) Rough coefficient cancellation.}
Retain the marked data and the hypotheses on \(F,\beta,H_m,H_n\), without
assuming a previously chosen \(\alpha\).  Fix \(0<\gamma<w_-<w_+<1\), and assume
\begin{equation}\label{eq:rough-dyad-range}
 [\log H_m/L,\log(2H_m)/L]\subseteq[w_-,w_+].
\end{equation}
For every interval \(I\subseteq[H_m,2H_m]\) and every \(|v|\le L^C\),
the coefficient
\begin{equation}\label{eq:untwisted-rough-coefficient}
 \alpha_m^{(0)}=m^{iv}\left(
   \one_{P^-(m)>x^\gamma}
   -\frac{D_\gamma(\log m/L)}{LV(W)}\one_{P^-(m)>W}\right)
 \quad(m\in I),
\end{equation}
set equal to zero outside \(I\), satisfies
\eqref{eq:marked-cancellation} with \(\alpha^{(0)}\) in place of
\(\alpha\), for every fixed \(A_0,B,A>0\).  This assertion is uniform in
the permitted dyads, interval, character, frequency \(t\), and \(v\).
Its constant and threshold may depend on the fixed parameters, including
\(\gamma,w_-,w_+\); in particular \(\gamma\) and \(w_--\gamma\) are
fixed before \(x\) grows.
\end{hypothesis}

Input~\ref{input:marked}(i) permits \(F\) to vary with \(x\), since the estimate is
uniform over all functions satisfying its two displayed conditions.
The next lemma inserts the fixed progression into the coefficient.

\begin{lemma}[A rough Type II estimate in a fixed progression]
\label{lem:twisted-rough-type-ii}
Fix \(\delta,C,D_*>0\) and \(0<\gamma<w_-<w_+<1\).  Let \(F,H_m,H_n,X\)
satisfy the hypotheses of Input~\ref{input:marked}, including
\eqref{eq:rough-dyad-range}.  Let \(\chi_0\) be any Dirichlet character
modulo the fixed \(M\).  On an arbitrary interval
\(I\subseteq[H_m,2H_m]\), put
\begin{equation}\label{eq:twisted-rough-coefficient}
 \alpha_m=\chi_0(m)m^{iv}\left(
   \one_{P^-(m)>x^\gamma}
   -\frac{D_\gamma(\log m/L)}{LV(W)}\one_{P^-(m)>W}\right),
 \qquad |v|\le L^C,
\end{equation}
and set \(\alpha_m=0\) outside \(I\).  Let \(\beta\) be supported on
\(W\)-rough integers in \([H_n,2H_n]\), with \(|\beta_n|\le L^C\).  Then
\begin{equation}\label{eq:twisted-rough-type-ii}
 \left|\sum_{m,n}\alpha_m\beta_nF(mn-1)\mathcal W(mn-1)\right|
      \ll XL^{-D_*}
\end{equation}
when \(K\) is sufficiently large in terms of the fixed data, independently
of the particular \(a_i\).  The estimate is uniform in the permitted
\(F\), dyads, intervals, \(\beta\), real \(v\), and characters \(\chi_0\).
The implicit constant and threshold may depend on all fixed data,
including \(M\) and the \(a_i\).
\end{lemma}

\begin{proof}
Let \(\alpha^{(0)}\) be the coefficient in
\eqref{eq:untwisted-rough-coefficient}.  For a character \(\chi\) modulo
\(k\le L^{A_0}\), the product \(\theta=\chi_0\chi\), with the stated
zero extensions, is a character modulo \(h=\operatorname{lcm}(M,k)\).
Indeed it is a homomorphism on the units modulo \(h\), and it vanishes
exactly on the nonunits modulo \(h\).  Since \(M\) is fixed,
\[
 h\le ML^{A_0}\le L^{A_0+1}
\]
for sufficiently large \(x\).  Input~\ref{input:marked}(ii), applied with
exponent \(A_0+1\), therefore gives
\[
 \left|\frac1{H_m}\sum_m\alpha_m\chi(m)m^{it}\right|
 =\left|\frac1{H_m}\sum_m\alpha_m^{(0)}\theta(m)m^{it}\right|
 \ll_A L^{-A}
 \qquad(|t|\le2XL^B).
\]
This includes imprimitive \(\theta\) and the case where \(\theta\) is
principal.  It is uniform in \(\chi_0\), since all the product characters
belong to the same enlarged modulus range.

These scalar sums and their permitted test ranges do not depend on \(K\)
or the \(a_i\): \(W\) is defined from \(x\) alone.  Fix their common
cancellation bounds using any one admissible auxiliary marked setup in
Input~\ref{input:marked}(ii).  The same scalar bounds then apply unchanged when
the \(K\) and bands used in Input~\ref{input:marked}(i) are chosen.

For large \(x\), \(x^\gamma>W\), so both terms of \(\alpha\) are supported
on \(W\)-rough integers.  The density is bounded on \([w_-,w_+]\), and
Mertens' theorem gives \(LV(W)\asymp L^{0.76}\).  Hence \(\alpha\) is
bounded, and in particular \(|\alpha_m|\le L^C\) eventually; multiplication
by \(\chi_0\) does not increase its size.  All the hypotheses of
Input~\ref{input:marked}(i) now hold, proving the lemma.
\end{proof}

The product progression will be imposed by the exact orthogonality identity
\begin{equation}\label{eq:product-progression}
 \one_{mn\equiv a\pmod M}
 =\frac1{\varphi(M)}\sum_{\chi_0\bmod M}
       \overline{\chi_0(a)}\chi_0(m)\chi_0(n).
\end{equation}
If either factor is a nonunit modulo \(M\), both sides are zero because
\(a\) is a unit.  For \(M=1\), the sum consists of the single trivial
character.  This identity leaves \(F\) with its required invariance.

We now choose the invariant function to encode the desired predecessors.
Fix a smooth function \(\Psi\) compactly supported in \((1,2)\), with
\(0\le\Psi\le1\) and \(A_\Psi=\int_1^2\Psi(u)\dd u>0\).  From now on take
\begin{equation}\label{eq:marked-family}
 \begin{split}
 F(h)&=\one_{\{h/h_{\mathcal P}=4Q\text{ for a prime }Q>x^{0.9}\}},
       \qquad h\ge1,\\
 w(d)&=\one_{d\equiv a\pmod M}\Psi(d/x)F(d-1)\mathcal W(d-1),
       \qquad d\ge2.
 \end{split}
\end{equation}
The quotient \(h/h_{\mathcal P}\) is unchanged by multiplication by any
group prime, even one already dividing \(h\), so this \(F\) is permitted
in Lemma~\ref{lem:twisted-rough-type-ii}.  For large \(x\), every group
prime is coprime to \(2M\) and lies in
\((\exp(L^{0.1}),\exp(L^{0.3}))\), while every prime \(Q>x^{0.9}\) lies
outside the groups.  Consequently, if \(w(d)>0\), there is a unique
representation
\begin{equation}\label{eq:marked-family-support}
 \begin{gathered}
 x<d<2x,\qquad d\equiv a\pmod M,\qquad d=4rQ+1,\\
 r=(d-1)_{\mathcal P}\in\mathcal G,\qquad Q>x^{0.9}\text{ prime}.
 \end{gathered}
\end{equation}
Moreover \(\mathcal W(d-1)=\mathcal W(r)\) and
\(r<(2x)/(4x^{0.9})\le x^{0.11}\).  We also have \(0\le w(d)\le B_K\).

For \(r\in\mathcal G\), the progression in \eqref{eq:marked-family}
prescribes the reduced class
\[
 b_r\equiv(a-1)(4r)^{-1}\pmod M
\]
for \(Q\).  It is reduced by \eqref{eq:admissible-progression}; inverses
and classes for \(M=1\) have their trivial meaning.  Define
\begin{equation}\label{eq:marked-mass-definitions}
 \begin{split}
 A_r&=\sum_{\substack{Q>x^{0.9}\text{ prime}\\Q\equiv b_r\pmod M}}
                  \Psi((4rQ+1)/x),\\
 X_0&=\sum_{d\ge2}w(d)=\sum_{r\in\mathcal G}\mathcal W(r)A_r,\\
 J_0&=\sum_{r\in\mathcal G}\frac{\mathcal W(r)}r.
 \end{split}
\end{equation}
The equality for \(X_0\) follows from the uniqueness in
\eqref{eq:marked-family-support}.  All these sums are nonnegative.

\begin{lemma}[Mass of the marked family]\label{lem:marked-mass}
For every fixed \(\vartheta>0\),
\begin{equation}\label{eq:marked-harmonic-bounds}
 1\le J_0\ll_K1,\qquad
 \sum_{\substack{r\in\mathcal G\\r>x^\vartheta}}
       \frac{\mathcal W(r)}r
       \ll_K\exp(-\vartheta L^{1-a_K}).
\end{equation}
For \(r\in\mathcal G\) with \(r\le x^{0.05}\), put
\[
 I_r=\frac{x}{4r}\int_1^2
       \frac{\Psi(u)}{\log((xu-1)/(4r))}\dd u.
\]
For every fixed \(A>0\), uniformly in these \(r\),
\begin{equation}\label{eq:progression-pnt}
 A_r=\frac{I_r}{\varphi(M)}+O_{M,\Psi,A}\left(\frac xr L^{-A}\right),
 \qquad I_r\asymp_\Psi\frac{x}{rL}.
\end{equation}
Furthermore,
\begin{equation}\label{eq:marked-mass-comparability}
 X_0\asymp_{M,\Psi}\frac{xJ_0}{L}.
\end{equation}
The comparison constants here can be chosen independently of \(K\) and
the \(a_i\); the threshold for \(x\) may depend on them.  For every fixed
\(\vartheta,A>0\),
\begin{equation}\label{eq:marked-tail}
 \sum_{\substack{r\in\mathcal G\\r>x^\vartheta}}
       \mathcal W(r)A_r=o(X_0L^{-A}).
\end{equation}
\end{lemma}

\begin{proof}
Put \(s=L^{-a_K}\).  Every group prime satisfies \(p^s\le e^2\), and for
large \(x\), \(p^{s-1}\le e^2/p\le1/2\).  The bound
\eqref{eq:mark-bound} and an Euler product including all prime powers give
\[
 \sum_{r\in\mathcal G}\mathcal W(r)r^{s-1}
 \le B_K\prod_{p\in\bigcup_i\mathcal P_i}(1-p^{s-1})^{-1}
 \le B_K\exp\left(2e^2\sum_iV_i\right)\ll_K1.
\]
This bounds \(J_0\) above.  Restricting \(J_0\) to
\(r=p_1\cdots p_K\), with one prime \(p_i\in\mathcal P_i\) to exponent
one, contributes exactly
\[
 \prod_{i=1}^K\left(\frac1{V_i}\sum_{p_i\in\mathcal P_i}\frac1{p_i}\right)
 =1.
\]
The disjoint groups make these products distinct.  Finally, Rankin's
inequality gives
\[
 \sum_{\substack{r\in\mathcal G\\r>x^\vartheta}}\frac{\mathcal W(r)}r
 \le x^{-\vartheta s}\sum_{r\in\mathcal G}\mathcal W(r)r^{s-1}
 \ll_K\exp(-\vartheta L^{1-a_K}),
\]
proving \eqref{eq:marked-harmonic-bounds}.

If \(r\le x^{0.05}\) and \(1\le u\le2\), then, for large \(x\),
\[
 \frac{x^{0.95}}8\le\frac{xu-1}{4r}\le\frac x2,
 \qquad
 0.9L\le\log\left(\frac{xu-1}{4r}\right)\le L.
\]
In particular the cutoff \(Q>x^{0.9}\) is automatic on the smooth
support.  The prime number theorem for the fixed modulus \(M\), uniformly
over its reduced classes and with arbitrary fixed logarithmic accuracy,
applied to \(t\mapsto\Psi((4rt+1)/x)\), proves the first assertion of
\eqref{eq:progression-pnt} by partial summation.  Indeed this test is
supported on \(t\asymp x/r\), whose logarithm is comparable to \(L\), and
its total variation is \(\int|\Psi'(u)|\dd u\), independently of \(r\).
The continuous prime main term is exactly \(I_r/\varphi(M)\) after the
substitution \(u=(4rt+1)/x\).  The displayed logarithmic bounds and
\(A_\Psi>0\) give the comparison for \(I_r\), with constants independent
of \(K\).

Write \(J_R=\sum_{r\in\mathcal G,\ r\le x^{0.05}}\mathcal W(r)/r\).
Taking \(A=3\) in \eqref{eq:progression-pnt} and summing with the marks
shows that \(\sum_{r\le x^{0.05}}\mathcal W(r)A_r\) is bounded above and
below by positive constants depending only on \(M,\Psi\) times
\(xJ_R/L\), once \(x\) is sufficiently large.  The summed error is
\(O_{M,\Psi}(xJ_RL^{-3})\).  For every \(r\), the smooth support forces
\(Q<x/(2r)\), so counting all positive integers instead of primes gives
\(0\le A_r\le x/(2r)\).  Equation~\eqref{eq:marked-harmonic-bounds}
therefore gives
\[
 J_0-J_R\ll_K e^{-0.05L^{1-a_K}},\qquad
 \sum_{\substack{r\in\mathcal G\\r>x^{0.05}}}\mathcal W(r)A_r
       \ll_K x e^{-0.05L^{1-a_K}}.
\]
Since \(J_0\ge1\), after a threshold depending on \(K,a_i\) we have
\(J_R\ge J_0/2\), and the last sum is at most \(xJ_0/L\).  Combining these
bounds proves \eqref{eq:marked-mass-comparability} with comparison
constants depending only on \(M,\Psi\).  For general fixed \(\vartheta>0\),
the same bound for \(A_r\) gives a tail
\(O_K(xe^{-\vartheta L^{1-a_K}})\).  Divide by
\(X_0L^{-A}\gg_{M,\Psi}xJ_0L^{-A-1}\) and use \(J_0\ge1\).  Since
\(L^{1-a_K}\) dominates every fixed multiple of \(\log L\), the quotient
tends to zero.  This proves \eqref{eq:marked-tail}.
\end{proof}

Fix \(0<\kappa<0.01\) and \(0<\epsilon<\kappa\), still before letting
\(x\) grow.  Lemma~\ref{lem:marked-mass} gives \(X_0\asymp xJ_0/L\) for the
total weight, with \(J_0\) the harmonic sum that combines the
estimates for fixed \(r\).  We now bound the summed divisibility remainders
for \(r\le x^\epsilon\); Equation~\eqref{eq:marked-tail} controls the
omitted tail.  For odd squarefree \(\ell\) coprime to \(M\), define
\begin{equation}\label{eq:distribution-data}
 \begin{split}
 D_0&=x^{1/2-\kappa/2},\\
 A_r(\ell)&=\sum_{\substack{Q>x^{0.9}\text{ prime}\\
                    Q\equiv b_r\pmod M\\\ell\mid4rQ+1}}
                    \Psi((4rQ+1)/x),\\
 g_r(\ell)&=\frac{\one_{\gcd(\ell,r)=1}}{\varphi(\ell)},
 \qquad R_r(\ell)=A_r(\ell)-g_r(\ell)A_r.
 \end{split}
\end{equation}
On these squarefree integers \(g_r\) is multiplicative, and
\(A_r(1)=A_r\), \(g_r(1)=1\), \(R_r(1)=0\).  If \(\gcd(\ell,r)>1\), both
terms in \(R_r(\ell)\) vanish.  Otherwise the new condition prescribes
the reduced class \(Q\equiv-(4r)^{-1}\pmod\ell\).  Together with
\(Q\equiv b_r\pmod M\), it gives one reduced class modulo \(M\ell\).

\begin{lemma}[Average distribution in the sieve moduli]
\label{lem:averaged-distribution}
For every fixed \(A>0\),
\begin{equation}\label{eq:averaged-distribution}
 \sum_{\substack{r\in\mathcal G\\r\le x^\epsilon}}\mathcal W(r)
 \sum_{\substack{\ell\le D_0\text{ odd and squarefree}\\
                  \gcd(\ell,M)=1}}
       |R_r(\ell)|\ll xJ_0L^{-A}.
\end{equation}
The constant and threshold may depend on the fixed parameters.
\end{lemma}

\begin{proof}
Fix \(r\le x^\epsilon\).  Since \(\epsilon<0.01<0.05\), the cutoff in
\(A_r(\ell)\) is automatic on its smooth support, as in
\eqref{eq:progression-pnt}.  Put \(\delta_0=(\kappa-\epsilon)/4>0\).
For every \(t\asymp x/r\) on that support, all the distinct moduli
\(M\ell\), \(\ell\le D_0\), satisfy
\[
 M\ell\le t^{1/2-\delta_0}
\]
for sufficiently large \(x\).  In fact, the exponent margin at the
smallest possible scale is
\[
 (1-\epsilon)(1/2-\delta_0)-(1/2-\kappa/2)
       =\frac{(\kappa-\epsilon)(1+\epsilon)}4>0,
\]
which absorbs the fixed factor \(M\) and the comparison constants.

The Bombieri--Vinogradov theorem \cite{Bombieri1965,Vinogradov1965}, restricted to this
fixed-power sublevel, gives for every fixed \(A'>0\) an
\(O(t(\log t)^{-A'})\) sum of absolute errors in prime counts up to \(t\),
using the maximum over reduced classes and main term
\(\operatorname{Li}(t)/\varphi(M\ell)\).  This follows from its usual
level \(t^{1/2}(\log t)^{-B}\), since \(\delta_0\) is fixed and positive.
Integrate those errors against the derivative of
\(\Psi((4rt+1)/x)\).  Its total absolute integral is bounded in terms of
\(\Psi\), and \(t\asymp x/r\), \(\log t\asymp L\) throughout.  The triangle
inequality under the integral thus gives
\[
 \sum_{\substack{\ell\le D_0\text{ odd and squarefree}\\
                  \gcd(\ell,M)=1}}
 \left|A_r(\ell)-g_r(\ell)\frac{I_r}{\varphi(M)}\right|
       \ll \frac xr L^{-A'}.
\]
For \(\gcd(\ell,r)=1\), the main term uses
\(\varphi(M\ell)=\varphi(M)\varphi(\ell)\); the other terms are zero.
The maximum over reduced classes in the theorem makes the bound uniform
in the class depending on \(r,\ell\).

To replace \(I_r/\varphi(M)\) by \(A_r\), observe that
\[
 \sum_{\substack{\ell\le D_0\\\ell\text{ squarefree}}}\frac1{\varphi(\ell)}
 \le\prod_{p\le D_0}\left(1+\frac1{p-1}\right)
 =V(D_0)^{-1}\ll L.
\]
Equation~\eqref{eq:progression-pnt}, with its logarithmic accuracy
increased by two, absorbs this factor.  Taking \(A'\) sufficiently large
therefore proves \(\sum_\ell|R_r(\ell)|\ll(x/r)L^{-A}\), uniformly in
the present \(r\).  Multiplying by \(\mathcal W(r)\), summing, and using
\(\sum_{r\le x^\epsilon}\mathcal W(r)/r\le J_0\) proves the lemma.
\end{proof}

The sieve products for this progression are
\begin{equation}\label{eq:progression-sieve-products}
 \begin{split}
 V_M(y)&=\prod_{\substack{2<p\le y\\p\nmid M}}
                  \left(1-\frac1{p-1}\right),\\
 \mathfrak S_M&=2\prod_{p>2}\left(1-\frac1{(p-1)^2}\right)
                \prod_{p\mid M}\left(1-\frac1{p-1}\right)^{-1}>0.
 \end{split}
\end{equation}
The infinite product is positive because its factors are positive and
their deviations from one have a convergent sum.  At every odd prime,
\[
 \frac{1-1/(p-1)}{1-1/p}=1-\frac1{(p-1)^2}.
\]
The factor at \(2\) and the omitted factors at primes dividing \(M\)
therefore give, with Mertens' theorem,
\begin{equation}\label{eq:progression-product-asymptotic}
 \frac{V_M(y)}{V(y)}\longrightarrow\mathfrak S_M,
 \qquad V(y)\sim\frac{e^{-\gamma_E}}{\log y}
 \quad(y\to\infty),
\end{equation}
where \(\gamma_E\) is Euler's constant.

Finally, for every \(r\in\mathcal G\) with \(r\le x^{0.11}\), its distinct
prime divisors exceed \(\exp(L^{0.1})\).  Their number is at most
\(0.11L^{0.9}\), and hence
\begin{equation}\label{eq:group-reciprocal-divisors}
 \sum_{p\mid r}\frac1p\le0.11L^{0.9}e^{-L^{0.1}}=o(1).
\end{equation}
For these primes, \(\log(1-1/(p-1))^{-1}\ll1/p\) with an absolute
constant.  It follows uniformly for all such \(r\) and all \(y\ge1\) that
\begin{equation}\label{eq:missing-group-factors}
 \begin{split}
 \prod_{\substack{2<p\le y\\p\nmid M}}(1-g_r(p))
 &=\prod_{\substack{2<p\le y\\p\nmid Mr}}
                \left(1-\frac1{p-1}\right)\\
 &=V_M(y)\left(1+O\bigl(L^{0.9}e^{-L^{0.1}}\bigr)\right)
  =V_M(y)(1+o(1)).
 \end{split}
\end{equation}
The constant in this error can be absolute; only the threshold uses the
fixed group exponents.
